\chapter{Microeconomía para directivos}\label{ch:microeconomics}

% Alcance: elasticidad, análisis marginal, estructura de mercado y teoría de subastas.
% Vínculos cruzados: elasticidad <-> respuesta de oferta y saturación en MMM (\cref{ch:paid-acquisition-platforms},
%   \cref{ch:attribution-and-measurement}); teoría de subastas <-> la subasta publicitaria (\cref{ch:paid-acquisition-platforms}).

Los mercados asignan recursos a través de precios; los directivos mueven precios sin comprender cómo
precio y cantidad interactúan en el margen. La elasticidad cuantifica esa interacción; el análisis
de estructura de mercado ubica a la empresa en el espectro de presión competitiva; la teoría de
subastas explica el mecanismo por el que dos de los mercados de medios más grandes de la historia
se liquidan en milisegundos. En conjunto, proporcionan una base microeconómica a cada decisión de
precio, posicionamiento y adquisición.

\section{Elasticidad del precio y análisis marginal}\label{sec:elasticity}
% complejidad: 3

Si se sube el precio, ¿se obtiene más o menos ingreso total? La respuesta depende enteramente de
cuán sensible es la demanda al precio --- pregunta que la elasticidad responde antes de comprometerse
con el cambio. Las curvas de demanda tienen pendiente negativa: mayor precio, menos unidades. Pero
el efecto sobre el ingreso de un cambio de precio depende de qué efecto domina --- el incremento
de precio por unidad o la pérdida de volumen. Un alza de precio genera más ingreso total únicamente
si los clientes son relativamente insensibles (demanda inelástica). En mercados elásticos, la
pérdida de volumen supera la ganancia de margen. El análisis marginal precisa esto: la regla de
maximización de utilidades es producir hasta que el ingreso marginal de la última unidad iguale su
costo marginal --- cualquier desviación deja dinero sobre la mesa.

La elasticidad precio de la demanda es el cambio porcentual en la cantidad demandada por cada uno
por ciento de cambio en el precio:
\begin{equation}\label{eq:elasticity}
  \varepsilon \;=\; \frac{\partial Q/Q}{\partial P/P}
              \;=\; \frac{\partial Q}{\partial P}\cdot\frac{P}{Q}.
\end{equation}
\(\varepsilon<0\) siempre (ley de la demanda). El punto de maximización del ingreso es \(\varepsilon=-1\)
(elasticidad unitaria); por debajo (inelástica, \(|\varepsilon|<1\)) un alza de precio eleva el ingreso;
por encima (elástica, \(|\varepsilon|>1\)) un alza de precio reduce el ingreso. El ingreso marginal se
vincula a la elasticidad mediante:
\[
  \mathrm{MR} \;=\; P\!\left(1+\frac{1}{\varepsilon}\right).
\]
Igualar \(\mathrm{MR}=\mathrm{MC}\) produce el máximo de utilidades estándar. El ingreso varía como:
\[
  \frac{\Delta R}{\Delta P} \;=\; Q\!\left(1 + \frac{1}{\varepsilon}\right),
\]
de modo que subir el precio en un mercado elástico (\(|\varepsilon|>1\)) genera un cambio negativo en el
ingreso --- la pérdida de volumen domina.

\Cref{eq:elasticity} supone una curva de demanda estática, preferencias constantes y sin respuesta
estratégica de los competidores. La elasticidad varía a lo largo de la curva de demanda (no es
constante para una función de demanda lineal), por lo que una calibración en un punto de precio no
se transfiere a un precio muy diferente. Las elasticidades de largo plazo son, en general, mayores
que las de corto plazo: los clientes necesitan tiempo para encontrar sustitutos, cambiar hábitos
o migrar de plataforma. En mercados digitales, el bloqueo de plataforma y los costos de cambio
pueden comprimir la elasticidad medida muy por debajo del valor estructural de largo plazo.

\begin{example}[La SaaS de referencia --- elasticidad del precio base]
La empresa SaaS de referencia fija su plan base en \money{50}/mes (\money{600}/año).
Supóngase que un estudio de demanda estima \(\varepsilon=-1.5\) al precio actual --- moderadamente elástica.
Un alza de precio propuesta del \qty{10}{\percent} a \money{55}/mes implica una caída de volumen de
\(1.5\times10\%=\qty{15}{\percent}\):
\[
  \Delta R \;\approx\; Q\!\left(1+\frac{1}{-1.5}\right)\times\Delta P
           = Q\times\frac{1}{3}\times\money{5}.
\]
El ingreso por cliente existente sube \money{5}, pero el \qty{15}{\percent} de los clientes genera
\emph{churn} --- para una base de \num{200} clientes, eso equivale a \num{30} suscripciones perdidas con
\(\mathrm{LTV}=\money{3150}\) cada una, un valor no generado de \money{94500}. Es poco probable que el
alza de precio sea positiva en términos netos a menos que la estimación de elasticidad esté materialmente
equivocada o exista un cohorte de clientes genuinamente inelástico. Conviene ejecutar el experimento
de precio en un segmento pequeño antes del lanzamiento general y medir la respuesta del churn
directamente, en lugar de depender del estimado agregado de \(\varepsilon\).
\end{example}

\begin{admonition}{Advertencia}
Usar la elasticidad agregada para una decisión de segmento oculta heterogeneidad peligrosa. Un
\(\varepsilon=-1.5\) global enmascara una distribución: algunos clientes son casi inelásticos
(empresas medianas y grandes, alto costo de cambio), otros muy elásticos (usuarios en prueba de
pymes). Un alza de precio uniforme no extrae valor de ninguno de los dos grupos de forma óptima.
La tarificación por segmento --- tema de \cref{ch:pricing} --- es la respuesta correcta; la
segmentación por elasticidad es su base empírica.
\end{admonition}

En el modelado de mezcla de \emph{marketing}, la curva de saturación --- rendimientos decrecientes
al gasto adicional --- es precisamente una elasticidad del lado de la demanda: la \emph{elasticidad
de respuesta} de las conversiones respecto a los dólares de medios, medida en el nivel de gasto
actual. A medida que el gasto aumenta, la respuesta marginal disminuye, razón por la que el
\(\mathrm{iROAS}\) (\cref{eq:iroas}) estimado por \textcite{hanssens2001marketing} y la
literatura de MMM es un estimado de elasticidad local alrededor del punto de presupuesto actual.
El modelo de adstock-saturación en \cref{ch:attribution-and-measurement} es la versión
temporalmente distribuida de \cref{eq:elasticity}: la misma disyuntiva entre gasto y rendimientos
decrecientes, extendida a lo largo de múltiples períodos mediante el decaimiento de adstock
(\cref{eq:adstock}).

La elasticidad alimenta directamente la arquitectura de precios (\cref{ch:pricing}), la
optimización del presupuesto de campañas (\cref{ch:paid-acquisition-platforms}) y las curvas de
saturación del modelado de mezcla de \emph{marketing} (\cref{ch:attribution-and-measurement}).
\textcite{pindyck2017} ofrece el tratamiento teórico completo; \textcite{mankiw2020} la
derivación de nivel introductorio.

\section{Estructura de mercado e índice de Lerner}\label{sec:lerner}
% complejidad: 3

¿Cuánto poder de fijación de precios tiene realmente la empresa? La estructura de mercado ---
el número y la simetría de los competidores --- determina el límite superior del margen sostenible.
El índice de Lerner convierte esa posición estructural en un número. La competencia perfecta fuerza
el precio al costo marginal; el monopolio puro permite el margen máximo. Entre esos extremos se
encuentran el oligopolio (pocas empresas, interdependencia estratégica) y la competencia
monopolística (muchas empresas, productos diferenciados). La mayoría de las empresas de tecnología
y SaaS operan en competencia monopolística: el poder de precio existe, pero está restringido por la
amenaza de sustitución. El índice de Lerner mide la fracción del precio que corresponde al margen,
vinculándose directamente con la elasticidad: un alto poder de precio requiere demanda inelástica.

El índice de Lerner es el margen precio-costo expresado como fracción del precio:
\begin{equation}\label{eq:lerner}
  L \;=\; \frac{P - \mathrm{MC}}{P} \;=\; -\frac{1}{\varepsilon}.
\end{equation}
La segunda igualdad se obtiene de igualar \(\mathrm{MR}=\mathrm{MC}\) en el máximo de utilidades y
sustituir \(\mathrm{MR}=P(1+1/\varepsilon)\). \(L\in[0,1]\): cero en competencia perfecta
(\(P=\mathrm{MC}\)), uno en monopolio puro (\(\mathrm{MC}=0\)). La fórmula codifica la misma
disyuntiva que \cref{eq:elasticity}: menor elasticidad (menor sensibilidad al precio) permite un
índice de Lerner más alto.

Tipología de estructura de mercado:
\begin{description}
  \item[Competencia perfecta] \(P=\mathrm{MC}\), \(L=0\). Muchas empresas idénticas, beneficio
    económico cero en el largo plazo. Agricultura de \emph{commodities}, algunos instrumentos
    financieros de mercado al contado.
  \item[Competencia monopolística] Productos diferenciados, entrada libre, \(L>0\) en el corto plazo
    pero erosionado por imitación. La mayoría de bienes de consumo, SaaS con muchos sustitutos
    cercanos.
  \item[Oligopolio] Pocas empresas grandes, interdependencia estratégica (el precio de cada empresa
    afecta a las rivales). Telecomunicaciones, infraestructura en la nube, publicidad en búsqueda.
  \item[Monopolio] Proveedor único, \(L\) máximo, sujeto a restricción regulatoria. Poco frecuente
    y regulado; un efecto de red fuerte (\cref{ch:growth}) puede aproximarlo temporalmente.
\end{description}

\Cref{eq:lerner} supone una curva de demanda estática y sin restricción regulatoria sobre precios.
En la práctica, índices de Lerner altos atraen escrutinio antimonopolio (la medida estándar en
economía de la competencia). La fórmula también asigna una única elasticidad a todo el mercado;
una empresa con múltiples segmentos enfrenta diferentes elasticidades --- y por lo tanto precios
óptimos distintos --- para cada uno. El equilibrio de largo plazo en competencia monopolística
lleva \(L\to 0\) a medida que los nuevos participantes imitan; el tiempo hasta el equilibrio depende
de qué tan defendible sea la diferenciación (tema de \cref{ch:advantage-and-disruption}).

\begin{example}[SaaS de referencia --- índice de Lerner]
La empresa SaaS enfrenta \(\varepsilon=-1.5\) al precio actual. Aplicando \cref{eq:lerner}:
\[
  L = -\frac{1}{-1.5} \approx 0.67.
\]
Un índice de Lerner de \num{0.67} significa que el \qty{67}{\percent} del precio del plan \money{50}
es margen sobre el costo marginal. Como referencia, el software de tipo \emph{commodity} se acerca
a \(L=0\) y el software empresarial frecuentemente supera \(L>0.80\); este resultado es consistente
con competencia monopolística con diferenciación moderada. La implicación gerencial: mayor
diferenciación --- integraciones, efectos de red, certificaciones --- puede reducir \(|\varepsilon|\)
(hacer la demanda menos elástica) y así ampliar el \(L\) sostenible sin cambios de precio.
\textcite{dranove7theconomics} presenta el rango empírico por industria.
\end{example}

\begin{admonition}{Advertencia}
Confundir el índice de Lerner con el margen bruto es un error frecuente y de consecuencias
importantes. El margen bruto es \((P-\mathrm{COGS})/P\); el índice de Lerner es
\((P-\mathrm{MC})/P\), donde \(\mathrm{MC}\) es el verdadero costo marginal --- el costo de
atender a un cliente adicional en el margen --- que en SaaS es cercano a cero. Un margen bruto del
\qty{70}{\percent} a escala no implica \(L=0.70\): los dos convergen sólo cuando no existen costos
fijos y el COGS captura todos los costos variables, lo cual no es cierto ni en la etapa de
crecimiento ni a escala.
\end{admonition}

En un oligopolio, el índice de Lerner se generaliza al modelo de Cournot: \(L_i = s_i / |\varepsilon|\),
donde \(s_i\) es la cuota de mercado de la empresa \(i\). Esto implica que las empresas con mayor
participación tienen más poder de precio incluso con la misma elasticidad de mercado. El Índice
Herfindahl-Hirschman (\cref{eq:hhi}) en \cref{ch:competitive-analysis} cuantifica la concentración;
el \(L\) promedio de la industria y el HHI están correlacionados en la mayoría de los estudios empíricos.

El razonamiento basado en el índice de Lerner se aplica directamente en la tarificación por valor
(\cref{ch:pricing}) e informa las decisiones de eficiencia de adquisición
(\cref{ch:paid-acquisition-platforms}). \textcite{pindyck2017} y \textcite{mankiw2020} derivan las
reglas de margen estándar; \textcite{dranove7theconomics} desarrolla las extensiones para oligopolio.

\section{Teoría de subastas y la subasta publicitaria}\label{sec:auction-theory}
% complejidad: 4 -> tabla es suficiente (caso discreto de 2 oferentes ilustra la estructura sin curvas)

Se participa en una subasta en tiempo real que se ejecuta en milisegundos. Conocer la estrategia
de oferta en equilibrio --- y por qué es óptimo ofertar el valor verdadero --- marca la diferencia
entre un gasto eficiente y el sobrepago sistemático. Una subasta de primer precio (se paga lo que
se oferta) incentiva a reducir la oferta por debajo del valor verdadero para evitar la maldición
del ganador; el descuento óptimo depende de las creencias sobre los valores de los competidores.
Una subasta de Vickrey (segundo precio) elimina ese incentivo: se paga únicamente la segunda oferta
más alta, por lo que ofertar el valor verdadero es una estrategia débilmente dominante
independientemente de lo que hagan los demás. Google y Meta ejecutan subastas de segundo precio
generalizado (GSP, por sus siglas en inglés) con un peso de calidad --- un multiplicador por
anunciante que convierte las ofertas brutas en ofertas efectivas, alterando el precio de cierre
y el ranking.

En la subasta de segundo precio (Vickrey) con \(n\) oferentes, la oferta veraz es una estrategia
dominante: si se gana (el valor verdadero \(v\) supera todos los valores rivales), se paga el
mayor valor rival \(v_{(2)}\) --- inferior a \(v\), lo que genera un excedente positivo. Si se
pierde, la oferta no tuvo efecto. Ninguna desviación de la oferta veraz mejora el resultado.

La subasta de segundo precio generalizado (GSP) con peso de calidad \(Q_i\) clasifica a los
oferentes por oferta efectiva \(b_i Q_i\) y cobra al ganador el monto necesario para superar al
siguiente anunciante clasificado:
\begin{equation}\label{eq:gsp}
  \mathrm{CPC}_i \;=\; \frac{b_{i+1}\,Q_{i+1}}{Q_i} + \delta,
\end{equation}
donde \(b_{i+1}Q_{i+1}\) es el AdRank efectivo del oferente clasificado justo por debajo de \(i\)
y \(\delta\) es un incremento pequeño (p.\ ej.\ \$0.01). Esta estructura --- ya derivada
operacionalmente en \cref{eq:ecpc} --- tiene una implicación clave: mejorar la calidad \(Q_i\)
reduce el costo por clic en cualquier ranking dado.

Ejemplo de dos oferentes con valores \(v_A=\money{5}\), \(v_B=\money{3}\) y puntajes de calidad
\(Q_A=10\), \(Q_B=8\):
\begin{center}\small
\begin{tabular}{lrrrl}
  \toprule
  Oferente & Valor real & Puntaje de calidad & Oferta efectiva & Resultado \\
  \midrule
  A        & \money{5}  & 10                 & 50              & Gana; paga \(3\times8/10+0.01=\money{2}.41\) \\
  B        & \money{3}  & 8                  & 24              & Pierde \\
  \bottomrule
\end{tabular}
\end{center}
El oferente A obtiene un excedente del consumidor de \(\money{5}-\money{2}.41=\money{2}.59\) por clic.
Si A redujera su calidad a \(Q_A=6\), la oferta efectiva caería a 30; la oferta efectiva de B de 24
sigue perdiendo, pero el CPC de A sube a \(3\times8/6+0.01=\money{4}.01\), reduciendo el excedente
a \money{0}.99. La mejora de calidad es inequívocamente rentable para el comprador.

\Cref{eq:gsp} supone que el puntaje de calidad es fijo y conocido al momento de la subasta. En
la práctica, tanto el Puntaje de Calidad de Google como la tasa de acción estimada de Meta
(\cref{eq:metavalue}) se actualizan continuamente; un creativo de alta calidad hoy puede ser
desplazado mañana por la predicción más baja de un modelo recién entrenado. El GSP no es
incentive-compatible en el mismo sentido preciso que Vickrey: los oferentes pueden beneficiarse
de reducir sus ofertas en entornos de múltiples posiciones \parencite{dranove7theconomics}. El
mecanismo también supone valores privados independientes; los valores correlacionados (comunes en
subastas de palabras clave de marca) abren espacio para la colusión o la coordinación tácita.

\begin{example}[SaaS de referencia --- optimización de la subasta GSP]
La empresa SaaS de referencia apunta a la palabra clave «software de gestión de proyectos» en
Google. Un competidor (B) con AdRank~16 ocupa la posición 2. Con Puntaje de Calidad \(Q_A=8\):
\[
  \mathrm{CPC} = 16/8 + \$0.01 = \money{2.01}.
\]
Con un presupuesto mensual de \money{120000} y un CPC de \money{2.01}, la empresa compra
\(\num{120000}/2.01\approx\num{59700}\) clics --- bien por encima del mínimo de la fase de
aprendizaje. Si se optimiza el creativo y la página de destino y se eleva \(Q_A\) a 10:
\[
  \mathrm{CPC} = 16/10 + \$0.01 = \money{1.61},\quad
  \text{clics} \approx \num{74500}\;\text{(+25\%)}.
\]
El mismo presupuesto compra un \qty{25}{\percent} más de tráfico a través de la calidad, no del
gasto. Esta es la forma operacional de \cref{eq:gsp} y es consistente con la derivación de CPC
en \cref{eq:ecpc}, \cref{ch:paid-acquisition-platforms}.
\end{example}

\begin{admonition}{Advertencia}
Ofertar por encima del valor verdadero para ganar la posición~1 en una subasta GSP reduce el
excedente por clic y puede destruir valor si el incremento de CTR de la posición~1 frente a la
posición~2 no compensa el CPC más alto. La decisión es un análisis marginal: ingreso incremental
por mejor posición frente a costo incremental --- exactamente la condición \(\mathrm{MR}=\mathrm{MC}\)
de \cref{eq:elasticity}.
\end{admonition}

El teorema de equivalencia de ingresos establece que, bajo condiciones estándar (oferentes neutros
al riesgo, valores privados independientes, distribuciones simétricas) todas las subastas estándar
--- de primer precio, segundo precio, inglesa, holandesa --- generan el mismo ingreso esperado para
el vendedor. El uso de pesos de calidad por parte de Google y Meta es una desviación de este caso
simétrico: permite a la plataforma extraer más ingresos clasificando a anunciantes de mayor calidad
y menores ofertas que generan más valor para el usuario (y por lo tanto más excedente del productor
del lado de la plataforma) de lo que el precio puro implicaría. La pregunta de diseño de mecanismos
--- qué subasta maximiza una suma ponderada del bienestar del anunciante y la plataforma --- es activa
en la literatura \parencite{mankiw2020}.

La estructura GSP de \cref{eq:gsp} es la base teórica de la mecánica de plataformas en
\cref{ch:paid-acquisition-platforms}; el puntaje de calidad es el componente de AdRank en
\cref{eq:adrank} y la fórmula de CPC en \cref{eq:ecpc}. \textcite{mankiw2020} y
\textcite{dranove7theconomics} cubren los fundamentos de la teoría de subastas; la implementación
en plataformas se desarrolla en \cref{ch:paid-acquisition-platforms}.

\section{Costos de transacción y fronteras de la empresa}\label{sec:transaction-costs}
% complejidad: 3

¿Debe la empresa construir su propio \emph{stack} de tecnología publicitaria, gestionar su propio
\emph{pipeline} de datos o usar las plataformas? La economía de costos de transacción otorga a la
decisión de fabricar-o-comprar un marco riguroso: internalizar cuando los mercados son demasiado
costosos de usar; externalizar cuando no lo son. Los mercados son mecanismos de asignación
eficientes, pero usarlos no es gratuito. Redactar y hacer cumplir contratos, monitorear el
desempeño de la contraparte y cambiar de proveedor consumen recursos. Cuando esos costos de
transacción superan los costos de gestionar una actividad dentro de la empresa, la internalización
es la elección eficiente. La intuición de Coase de 1937 es que las fronteras de la empresa se
fijan en el margen donde el costo de una transacción interna adicional iguala el costo de una
transacción de mercado adicional.

La decisión de fabricar-o-comprar se evalúa en cuatro dimensiones:
\begin{description}
  \item[Especificidad de activos] ¿Qué tan especializado es el activo para esta relación? Alta
    especificidad (integraciones personalizadas, datos propietarios) crea riesgo de \emph{hold-up}
    en una relación de mercado --- la contraparte puede extraer rentas en la renegociación.
    Internalizar.
  \item[Incertidumbre] Los entornos más difíciles de contratar (tecnología que cambia rápidamente,
    volúmenes impredecibles) elevan el costo de los contratos de mercado. Internalizar o usar
    contratación relacional (de largo plazo, basada en confianza).
  \item[Frecuencia] Las transacciones frecuentes y estandarizadas reducen el costo por unidad de
    la contratación de mercado en relación con el costo fijo del gobierno interno. Externalizar.
  \item[Dificultad de medición] Cuando la calidad del resultado es difícil de observar y verificar,
    los contratos de mercado generan riesgo moral (la contraparte no cumple porque la baja calidad
    no se penaliza). Internalizar o invertir en monitoreo.
\end{description}
\parencite{coase1937nature}.

El \emph{framework} supone que la internalización resuelve efectivamente el problema de gobierno.
Las grandes empresas también enfrentan costos de transacción internos --- burocracia, incentivos
desalineados, fallas de coordinación --- que pueden superar los costos de mercado a escala. La
frontera óptima no es, por tanto, «empresa grande vs.\ mercado», sino una comparación continua de
costos en cada margen del alcance organizacional. Los ecosistemas de plataforma crean un tercer
modo: mercados gobernados (App Store, API de Google Ads) que combinan precios de mercado con
estándares impuestos por la plataforma, internalizando parcialmente el riesgo de especificidad
de activos sin integración vertical completa \parencite{dranove7theconomics}.

\begin{example}[SaaS de referencia --- \emph{stack} de atribución propio vs. plataformas]
La empresa SaaS evalúa si construir un \emph{stack} de atribución propia o depender de las APIs
nativas de conversión de las plataformas. Aplicando el \emph{framework}:
\begin{itemize}
  \item \emph{Especificidad de activos}: los datos de conversión propios son altamente específicos
    (identificadores de clientes propietarios, eventos del lado del servidor); alto riesgo de
    \emph{hold-up} si se almacenan exclusivamente en una herramienta de terceros. Señal: construir.
  \item \emph{Incertidumbre}: las regulaciones de privacidad (web sin \emph{cookies}) cambian
    rápido; un \emph{stack} personalizado requiere inversión continua de ingeniería. Señal: comprar
    o usar la API de la plataforma (menor mantenimiento).
  \item \emph{Frecuencia}: los eventos de conversión ocurren continuamente a escala; las APIs
    estandarizadas los manejan a bajo costo. Señal: comprar.
  \item \emph{Dificultad de medición}: las conversiones reportadas por la plataforma pueden
    sobrecontarse debido a las ventanas de atribución; la verificación requiere una capa de medición
    independiente. Señal: construir la capa de verificación; externalizar la entrega.
\end{itemize}
Decisión: usar la API de Conversiones de la plataforma (\cref{ch:attribution-and-measurement}) para
la entrega, construir una capa de datos propia para la verificación y el modelado. El enfoque
híbrido minimiza tanto los costos de transacción como el riesgo de \emph{hold-up}.
\end{example}

\begin{admonition}{Advertencia}
Las decisiones de fabricar-o-comprar tomadas sobre costos hundidos en lugar de costos marginales
futuros están sistemáticamente sesgadas. Una empresa que ya ha invertido en un \emph{stack}
propietario subestimará el costo de mantenimiento continuo y sobreestimará el costo de cambio,
ambos factores favorecen el \emph{statu quo} independientemente de la economía actual. La aversión
a las pérdidas (\cref{eq:prospect}) amplifica esto: la perspectiva de «perder» la inversión existente
se codifica como una pérdida cierta, ponderada en \(\lambda\approx2.25\).
\end{admonition}

La integración vertical y el gobierno de plataformas son extensiones del \emph{framework} de Coase
aplicadas a los ecosistemas digitales. Una plataforma que posee tanto el mercado como un producto
competidor enfrenta el mismo problema de fabricar-o-comprar que un comprador, pero invertido:
internaliza cuando puede ganar más del autoabastecimiento que de la externalidad que genera al
habilitar la competencia de terceros. Esta es la economía detrás de las disputas antimonopolio
en las tiendas de aplicaciones: ¿la integración vertical de la plataforma (cobrar una comisión,
favorecer sus propios productos) es una reducción de costos de transacción que mejora el bienestar
o es un cierre anticompetitivo? \textcite{coase1937nature} proporciona el \emph{framework}
fundacional; \textcite{dranove7theconomics} lo aplica a los mercados de plataformas.

El análisis de fabricar-o-comprar reaparece en la estrategia de operaciones (\cref{ch:operations})
y en la discusión de las decisiones de \emph{stack} de tecnología de \emph{marketing} en
\cref{ch:attribution-and-measurement}. \textcite{coase1937nature} es el texto fundacional;
\textcite{dranove7theconomics} lo extiende a los mercados digitales.
